\section{Заключение}
Исследована задача выбора способа обучения двухмодельного
маршрутизатора по объёму и составу парной разметки и структуре
ошибок экспертов. Полный контроль включает восемь обучаемых методов,
два необучаемых правила, 21 зависимую пару и 110\,880 обученных роутеров.
Качество каскада, профили спасений/вреда и качество выбора по dev
рассматриваются как разные объекты анализа.

На natural попарные accuracy-различия не подтверждены общей поправкой.
При harm-обогащении обнаружено 56 взаимодействий метода и состава:
семь gain-оценивателей чувствительнее к смещению цели, чем weak-error
и confidence, на четырёх исходных парах. Между самими семью
gain-параметризациями взаимодействия не подтверждены; четыре простых
различия в harm8 относятся к одной low-label паре. Надёжная смена
знака преимущества, поддержанная интервалами в обоих режимах,
не установлена. Численные различия natural-профилей остаются
описательными после общей поправки и не доказывают равенства accuracy.

Следовательно, гипотеза о неодинаковой реакции методов на дисбаланс
поддержана для сравнения семейств обучающих целей, но пока не даёт
универсальной карты выбора внутри gain-семейства. Ухудшение после
изменения состава не является улучшением обучения: в основной семье
не обнаружено положительных эффектов raw-обогащения над natural.
IPW-контроль показывает роль смещения обучающей цели; он не заменяет
задачу выбора маршрутизатора. Выбор по отдельному dev также требует
проверки, а не назначения лидера по тестовым средним.

Следующий этап --- независимые датасеты и заново обученные эксперты,
проверка заранее заданного правила выбора на новых задачах,
согласование класса функций и tuning-критерия. Дополнительные
повторения на том же фиксированном тесте не заменяют эту проверку.
